\section{Experimental Evaluation}
\label{sec:experiments}

In this section, we present the empirical evaluation of \textbf{QWS-Merge} against competitive baselines, analyzing performance under both homogeneous and heterogeneous evaluation streams, and investigating sensitivity to test batch sizes.

\subsection{Experimental Setup}
\label{subsec:setup}
To evaluate model merging approaches in the presence of severe task conflict and restricted parameter capacity, we construct a heterogeneous multi-task benchmark spanning four distinct visual domains:
\begin{enumerate}
    \item \textbf{MNIST:} Handwritten digits ($10$ classes).
    \item \textbf{FashionMNIST:} Fashion products ($10$ classes).
    \item \textbf{CIFAR-10:} Natural objects ($10$ classes).
    \item \textbf{SVHN:} Street View House Numbers ($10$ classes).
\end{enumerate}

\paragraph{Backbone Model:} We employ a compact Vision Transformer backbone, specifically $\mathtt{vit\_tiny\_patch16\_224}$ \cite{dosovitskiy2020image}, containing only $5.7$M parameters. The model comprises $14$ layer groups: a Patch Embedding layer, $12$ Transformer blocks, and a final Layer Normalization layer. This restricted capacity makes the model highly susceptible to parameter-space interference.

\paragraph{Converged Expert Training:} To address the ``fake expert'' limitation identified in prior drafts, each expert is trained to high convergence. We use a dual learning rate schedule ($2 \times 10^{-5}$ for the ViT backbone and $10^{-3}$ for the classification heads) and a training batch size of $64$, training for $15$ epochs using AdamW. This converges to a high-performance specialized ceiling across all four tasks, as shown in Table~\ref{tab:main_results}.

\paragraph{Few-Shot Calibration Set:} For methods requiring parameter tuning (OFS-Tune, Linear Router, and QWS-Merge), we use a tiny offline calibration validation set containing only $16$ samples per task ($64$ total calibration images).

\subsection{Baselines}
We compare QWS-Merge against five competitive baselines:
\begin{itemize}
    \item \textbf{Individual Experts (Ceiling):} Specialized, task-specific networks. This represents the empirical upper bound of performance.
    \item \textbf{Uniform Merging (Task Arithmetic):} Static, uniform merging of task vectors \cite{ilharco2022editing} with a scaling coefficient of $\lambda = 0.3$.
    \item \textbf{AdaMerging (Unsupervised TTA):} Unsupervised test-time adaptation \cite{liang2023adamerging}, optimizing $56$ layer-wise coefficients to minimize test stream prediction entropy.
    \item \textbf{OFS-Tune (Supervised Static):} Supervised static layer-wise coefficient optimization on the $64$-sample calibration set via Adam.
    \item \textbf{Linear Router (Classical Baseline):} Our proposed baseline, mapping the input's pooled backbone representations directly to routing weights via a soft linear layer, trained on the $64$-sample calibration set for $100$ steps.
\end{itemize}

\subsection{Quantitative Results}
\label{subsec:results}

\paragraph{Homogeneous Stream Performance:}
The performance of all methods on a homogeneous test stream (where batches contain samples from a single task) is summarized in Table~\ref{tab:main_results}.

\begin{table*}[t]
\caption{Multi-task image classification accuracy ($\%$) on the homogeneous $\mathtt{vit\_tiny\_patch16\_224}$ benchmark. Best multi-task results (excluding the individual expert ceiling) are highlighted in \textbf{bold}.}
\label{tab:main_results}
\vskip 0.15in
\begin{center}
\begin{small}
\begin{sc}
\begin{tabular}{lcccccr}
\toprule
Method & MNIST & FashionMNIST & CIFAR10 & SVHN & Joint Mean \\
\midrule
Individual Experts (Ceiling)   & 92.50\% & 77.70\% & 77.40\% & 34.50\% & 70.52\% \\
\midrule
Uniform Merge (TA, $\lambda = 0.3$) & 63.80\%  & 44.50\%  & 67.20\% & 21.90\% & 49.35\% \\
AdaMerging (Unsupervised TTA) & 78.00\% & 52.80\%  & 73.70\% & 23.80\% & 57.07\%  \\
OFS-Tune (Supervised Static)  & 67.20\%  & 58.40\% & 67.40\% & 27.00\% & 55.00\% \\
Linear Router (Classical Baseline) & \textbf{91.20\%} & \textbf{67.00\%} & \textbf{71.40\%} & 15.30\% & \textbf{61.23\%} \\
\midrule
\textbf{QWS-Merge (Ours)}      & 77.60\% & 63.50\% & 64.60\% & \textbf{31.60\%} & 59.32\% \\
\bottomrule
\end{tabular}
\end{sc}
\end{small}
\end{center}
\vskip -0.1in
\end{table*}

\paragraph{Heterogeneous Stream Performance:}
To evaluate the impact of task mixing and batch sizes, we feed the models a randomly shuffled, heterogeneous multi-task test stream across batch sizes $B \in \{1, 16, 256\}$. The comparative results are detailed in Table~\ref{tab:hetero_results}.

\begin{table}[h]
\caption{Multi-task accuracy ($\%$) under heterogeneous (mixed) test streams across different batch sizes ($B$).}
\label{tab:hetero_results}
\vskip 0.15in
\begin{center}
\begin{small}
\begin{sc}
\begin{tabular}{lccc}
\toprule
Method & $B=1$ & $B=16$ & $B=256$ \\
\midrule
Uniform Merge & 49.20\% & 49.20\% & 49.20\% \\
AdaMerging    & 57.20\% & 57.20\% & 57.20\% \\
OFS-Tune      & 55.60\% & 55.60\% & 55.60\% \\
\midrule
Linear Router & \textbf{55.70\%} & 48.60\% & 47.70\% \\
\textbf{QWS-Merge (Ours)} & 54.90\% & \textbf{48.80\%} & \textbf{48.70\%} \\
\bottomrule
\end{tabular}
\end{sc}
\end{small}
\end{center}
\vskip -0.1in
\end{table}

\subsection{Deep Empirical Analysis}
\label{subsec:analysis}

\paragraph{Catastrophic Representational Collapse:}
When individual experts are trained to true convergence, standard static merging methods suffer severely from parameter conflict. Uniform merging collapses to $49.35\%$ joint mean accuracy, representing a loss of over $21\%$ in absolute accuracy compared to the expert ceiling. This performance drop is particularly severe on FashionMNIST ($44.50\%$ vs $77.70\%$) and SVHN ($21.90\%$ vs $34.50\%$). Because the ViT-Tiny's parameter space is extremely compact, static blending forces conflicting task-specific features to cancel each other out, resulting in destructive representational interference.

\paragraph{The Double-Edged Sword of Linear Routing:}
The classical Linear Router baseline performs strongly on low-conflict datasets like MNIST ($91.20\%$) and FashionMNIST ($67.00\%$), achieving a joint mean of $61.23\%$. However, because its routing weights are unconstrained and optimized globally, the Linear Router is highly vulnerable to catastrophic collapse under extreme task conflict. This is vividly illustrated on the SVHN task (which has the largest domain shift), where the Linear Router collapses to a near-random $15.30\%$ accuracy, losing more than half of the expert capacity.

\paragraph{Wave-Like Subspace Regularization in QWS-Merge:}
In contrast to the classical router, our proposed QWS-Merge exhibits extraordinary regularization properties. On the highly challenging SVHN dataset, QWS-Merge maintains an accuracy of \textbf{31.60\%}, preserving $91.5\%$ of the specialized expert ceiling and outperforming the Linear Router by an absolute margin of \textbf{$+16.30\%$}. This confirms that the non-monotonic cosine wave phase modulation and spherical projections act as a highly regularized, low-noise coordination subspace, preventing the unconstrained parameter-space collapse and overfitting that afflict classical linear routing in high-conflict, data-scarce regimes.

\begin{figure}[ht]
\vskip 0.1in
\begin{center}
\centerline{\includegraphics[width=0.9\columnwidth]{comparison_plot.png}}
\caption{Performance comparison ($\%$) under homogeneous streams across the multi-task benchmark.}
\label{fig:comparison_plot}
\end{center}
\vskip -0.1in
\end{figure}

\begin{figure}[ht]
\vskip 0.1in
\begin{center}
\centerline{\includegraphics[width=0.9\columnwidth]{heterogeneous_plot.png}}
\caption{The impact of batch size ($B$) on multi-task accuracy under a mixed-task heterogeneous test stream.}
\label{fig:heterogeneous_plot}
\end{center}
\vskip -0.1in
\end{figure}

\paragraph{Batch Heterogeneity \& "Heterogeneity Collapse":}
A core, scientifically transparent contribution of our work is the study of batch size and task mixing on dynamic merging (Table~\ref{tab:hetero_results} and Figure~\ref{fig:heterogeneous_plot}):
\begin{enumerate}
    \item \textbf{Static Methods:} Methods like Uniform, AdaMerging, and OFS-Tune are batch-invariant (constant across $B=1, 16, 256$) because their coefficients are frozen.
    \item \textbf{Dynamic Methods:} Both the Linear Router and QWS-Merge exhibit clear batch-size sensitivity. At $B=1$ (fully sample-wise dynamic), they perform at their best because there is no task mixing. However, at $B=256$, task representations are mixed and averaged across the batch during the wavefunction collapse step, causing the dynamic coefficients to collapse back toward uniform compromises (the Linear Router drops to $47.70\%$, and QWS-Merge drops to $48.70\%$, performing below static AdaMerging).
\end{enumerate}
Crucially, QWS-Merge shows superior resilience to this heterogeneity collapse than the Linear Router: at $B=256$, QWS-Merge is $1.0\%$ more accurate than the Linear Router. This highlights that wave-inspired phase-interference maintains a more robust parameter subspace even under high batch-level task mixing.

\subsection{Limitations and Deployment Discussion}
\label{subsec:limitations}

While QWS-Merge presents a compelling, regularized approach to dynamic model merging, we transparently address several limitations that affect its current formulation and restrict its immediate real-world deployability.

\paragraph{Batch Dependency and I.I.D. Violation:}
A major theoretical limitation of QWS-Merge is its batch-dependent inference. Because sample-level routing coefficients $\alpha_{k, b}(l)$ are averaged across the batch dimension to compute a batch-level coefficient $\bar{\alpha}_k(l)$, the parameters used to process an individual sample $x_b$ depend on the other samples in the same batch. This violates the independent-and-identically-distributed (I.I.D.) assumption of standard machine learning. Consequently, evaluating the same image in different batches (e.g., in a pure MNIST batch vs a mixed-task batch) yields slightly different predictions. 

To mitigate this and make the model deployable in standard online single-sample inference streams ($B=1$), future work could employ an Exponential Moving Average (EMA) or a rolling queue of recent routing coefficients during inference, or map sample-level routing directly to localized activation-routing layers (reminiscent of Mixture-of-Experts) rather than dynamically reconstructing the entire backbone weight matrix.

\paragraph{Capacity-Regularization Trade-Off:}
Our results highlight a clear capacity-regularization trade-off between the classical Linear Router and QWS-Merge. The Linear Router achieves a slightly higher overall joint mean accuracy ($61.23\%$) by exploiting its unconstrained linear projection matrix, which preserves high performance on low-conflict tasks (MNIST, FashionMNIST). However, this unconstrained parameter space collapses under extreme conflict (SVHN: $15.30\%$). QWS-Merge, by bounding the parameters inside a layer-wise cosine wave-interference subspace, acts as a heavy regularizer, sacrificing a small amount of capacity on simple tasks to achieve exceptional robustness under extreme conflicts (SVHN: $31.60\%$).

\paragraph{Few-Shot Calibration and Overfitting Risks:}
Both the Linear Router and QWS-Merge calibrate their routing parameters on a tiny validation set of $64$ samples. Optimizing these parameters for $100$ steps carries a risk of overfitting to the validation split. While QWS-Merge's ultra-compact parameter footprint ($336$ parameters) mitigates this risk far better than the Linear Router ($772$ parameters), evaluating on larger, out-of-distribution calibration sets remains a key direction to confirm the generalization bounds of dynamic parameter-space merging.
